\documentclass[11pt]{article}

\usepackage[margin=1in]{geometry}
\usepackage{fancyhdr}

\title{\Large PLUS Session Outline \\ \normalsize \textit{MATH 1554:} Linear Algebra}
\author{Ahmed Mohammed}
\date{27th of June, 2023}

% ------- Set up header -------
\renewcommand{\sectionmark}[1]{\markright{#1}} % from https://tex.stackexchange.com/a/317393/227665
\pagestyle{fancy}
\fancyhead[L]{27th of June, 2023}
% \fancyhead[C]{\textit{\nouppercase{\rightmark}}}
\fancyhead[C]{\textit{MATH 1554 PLUS Session Outline}}
\fancyhead[R]{Mohammed \textbf{\thepage}}
% -------------------------------------------------------

\begin{document}
%\maketitle
\thispagestyle{fancy}
\section{\centering Open Activity}
	\begin{enumerate}
		\item \textit{Activity}. The group(s) will be asked a conceptual question and will be given a moment to answer it. Then, the group will be introduced to the 
			main topic of the session.
		\item \textit{Collaboration}. The group(s) will be expected to ponder upon the question together in order to reach to an answer.
		\item \textit{Intended Goal}. The intention of this activity is to serve as a motivation for the topic of this session. It (partially) answers \textit{why} we
			care about it.
	\end{enumerate}

\section{\centering Main Activity}
	\begin{enumerate}
		\item \textit{Activity}. The answer to the previous conceptual problem will be revealed in the form of another question which I'll go through. Students will
			then be asked to answer more questions one by one, with me providing them with enough time for them to answer. I'll then go through each 
			question.
		\item \textit{Collaboration}. During the amount of time the students are pondering upon each question, they are expected to work with each other.
		\item \textit{Intended Goal}. This activity furthermore shows \textit{why} the main topic of the lecture is powerful. It also serves as a practice for 
			answering similar questions found in previous exams.
	\end{enumerate}

\section{\centering Closing Activity}
	\begin{enumerate}
		\item \textit{Activity}. In this activity, students will be lead to other resources that discuss the main topic of this session in more detail, or in an alternative
			fashion. Then, the next session's main topic is ``foreshadowed."
		\item \textit{Collaboration}. \textsc{N/A}
		\item \textit{Intended Goal}. This activity provides students with more resources, and introduces the topic of the next session.
	\end{enumerate}
\end{document}